\documentclass[final,table]{beamer} % 【修复2】加入 table 选项支持 \rowcolor
\usepackage[T1]{fontenc}
\usepackage[utf8]{inputenc}
\usepackage[orientation=portrait,size=a0,scale=1.4]{beamerposter}
\usepackage{graphicx}
\usepackage{booktabs}
\usepackage{pifont}
\usepackage{tikz}
\usetikzlibrary{shapes,arrows,positioning,shadows,calc,backgrounds}
\usepackage{xcolor}
\usepackage[skins,most]{tcolorbox}
% 【修复3】移除了容易与 beamer 冲突的 enumitem 包

% --- Color Palette (Matched to Reference Image) ---
\definecolor{bgcanvas}{HTML}{F6F6F6}     
\definecolor{titlebg}{HTML}{EBE6E1}      
\definecolor{boxbg}{HTML}{F0F4F8}        
\definecolor{textdark}{HTML}{333333}     
\definecolor{accentmain}{HTML}{667A8A}   
\definecolor{accentline}{HTML}{BFC9D4}   
\definecolor{numberbg}{HTML}{FFFFFF}     
\definecolor{chartline}{HTML}{506780}    

\setbeamercolor{background canvas}{bg=bgcanvas}
\color{textdark} 

% Beamer Native Itemize Customization
\setbeamertemplate{itemize item}{\color{accentmain}\Large\textbullet}
\setbeamertemplate{itemize subitem}{\color{accentmain}\large\textbullet}

% --- Custom Title Box with Logo Placeholder ---
\newcommand{\maintitleblock}{
    \begin{tikzpicture}[remember picture, overlay]
        \node[anchor=north, yshift=-4cm] at (current page.north) {
            \begin{tcolorbox}[
                colback=titlebg,
                colframe=titlebg,
                width=0.94\paperwidth,
                height=22cm,
                arc=30pt,
                boxrule=0pt,
                halign=center,
                valign=center,
                drop fuzzy shadow
            ]
                % Logo Area (Left)
                \begin{tikzpicture}[remember picture, overlay]
                    \node[anchor=west, xshift=2cm] at (frame.west) {
                        \begin{tikzpicture}
                            \draw[accentmain, line width=4pt] (0,0) circle (3cm);
                            \node at (0,0) {\sffamily \Huge \textbf{LOGO}};
                            \node[below=3.5cm] at (0,0) {\sffamily \Large \textbf{LOREM IPSUM}};
                            \node[below=4.5cm] at (0,0) {\sffamily \normalsize SINCE 1911};
                        \end{tikzpicture}
                    };
                \end{tikzpicture}
                
                % Title Content
                \begin{minipage}{0.65\paperwidth}
                    \centering
                    {\fontsize{85}{100}\selectfont \color{textdark} \textbf{Loyalty Signaling and Discursive Convergence:} \\[0.8cm]}
                    {\fontsize{60}{70}\selectfont \color{textdark!80} \textit{How Hong Kong Chief Executives Use Beijing Rhetoric} \\ \textit{to Signal Political Commitment} \\[2cm]}
                    {\fontsize{45}{55}\selectfont \color{textdark} \textbf{Author Name$^{1}$}, \textbf{Co-Author Name$^{2}$} \\[0.8cm]}
                    {\fontsize{32}{40}\selectfont \color{textdark!70} $^1$ Department of Political Science, University of Hong Kong \\ $^2$ Institute of Public Policy \par}
                \end{minipage}
            \end{tcolorbox}
        };
    \end{tikzpicture}
}

% --- Custom Numbered Header Style ---
% 【修复1】这里改为了 .style 2 args={
\tcbset{
    posterbox/.style 2 args={
        enhanced,
        colback=boxbg,
        colframe=boxbg,
        boxrule=0pt,
        arc=20pt,
        left=2cm, right=2cm, top=2cm, bottom=2cm,
        fontupper=\Large\color{textdark},
        title={}, 
        overlay={
            \fill[white, rounded corners=15pt] ([xshift=1.5cm,yshift=1.5cm]frame.south west) rectangle ([xshift=-1.5cm,yshift=-1.5cm]frame.north east);
            \node[circle, fill=numberbg, draw=accentmain!30, line width=2pt, minimum size=2.5cm, font=\fontsize{50}{60}\selectfont\bfseries\color{accentmain}, anchor=west, drop shadow={opacity=0.1}] 
            at ([xshift=1cm, yshift=-3cm]frame.north west) {#1};
            \node[anchor=west, xshift=4cm, font=\fontsize{55}{65}\selectfont\bfseries\color{textdark}] 
            at ([yshift=-3cm]frame.north west) {#2};
            \draw[accentline, line width=3pt] ([xshift=1cm, yshift=-5.5cm]frame.north west) -- ([xshift=-1cm, yshift=-5.5cm]frame.north east);
        },
        top=6.5cm 
    }
}

\begin{document}
\begin{frame}[t]

% Insert the top title block
\maintitleblock

% Push main content down below the title box
\vspace*{30cm} 

% --- MAIN CONTENT GRID (2 Columns) ---
\begin{columns}[T]

    % ================= LEFT COLUMN =================
    \begin{column}{0.45\textwidth}
        
        % --- 1. INTRODUCTION ---
        \begin{tcolorbox}[posterbox={1}{Introduction}]
            \textbf{Research Question:} Why do Hong Kong Chief Executives (CEs) increasingly adopt mainland Chinese political terminology?
            
            \vspace{1.5cm}
            \textbf{Key Phenomenon:}
            \begin{itemize}
                \item \textit{Mainlandization} (mainland-ization) of HK executive discourse.
                \item Growing use of \textit{guanqiang} (官腔) — bureaucratic/state-speak.
                \item Keywords: 爱国者治港 (Patriots governing HK), 新质生产力 (New Quality Productive Forces), 大湾区 (Greater Bay Area) integration.
            \end{itemize}
            
            \vspace{1.5cm}
            \begin{tcolorbox}[colback=accentmain!10, colframe=accentmain!30, boxrule=2pt, arc=15pt, left=1cm, right=1cm, top=1cm, bottom=1cm]
                \textbf{Argument:} Rhetoric convergence acts as a strategic \textbf{loyalty signal} to Beijing following political crises (e.g., 2019 protests, regulatory reforms), overriding local dissent to secure central backing.
            \end{tcolorbox}
        \end{tcolorbox}

        \vspace{2cm}

        % --- 3. RESULTS 1: TRENDS ---
        \begin{tcolorbox}[posterbox={3}{Results 1: Loyalty Trends}]
            Officialese frequency analysis across five CE administrations showing exponential growth.\\[1.5cm]
            
            \centering
            \begin{tikzpicture}[x=10cm, y=0.15cm]
                \draw[black!10, dashed, line width=3pt] (1,0) grid[ystep=20, xstep=0.5] (3,80);
                \draw[-, line width=5pt, textdark!40] (0.8,0) -- (3.2,0);
                \draw[-, line width=5pt, textdark!40] (0.8,-5) -- (0.8,85);
                
                \node[left, font=\LARGE, text=textdark!80] at (0.6, 0) {0};
                \node[left, font=\LARGE, text=textdark!80] at (0.6, 40) {40};
                \node[left, font=\LARGE, text=textdark!80] at (0.6, 80) {80};
                \node[rotate=90, font=\LARGE, yshift=2.5cm] at (0.8, 40) {Density (Hits/10k)};
                
                \node[below, font=\LARGE, align=center, yshift=-0.5cm] at (1,-5) {Tung \\ (\small 1997-2005)};
                \node[below, font=\LARGE, align=center, yshift=-0.5cm] at (1.5,-5) {Tsang \\ (\small 2005-2012)};
                \node[below, font=\LARGE, align=center, yshift=-0.5cm] at (2,-5) {Leung \\ (\small 2012-2017)};
                \node[below, font=\LARGE, align=center, yshift=-0.5cm] at (2.5,-5) {Lam \\ (\small 2017-2022)};
                \node[below, font=\LARGE, align=center, yshift=-0.5cm] at (3,-5) {Lee \\ (\small 2022-2024)};

                \draw[color=chartline, line width=8pt, rounded corners=15pt, drop shadow] 
                    plot coordinates {(1,12.3) (1.5,18.7) (2,25.1) (2.5,41.8) (3,68.4)};
                
                \foreach \x/\y in {1/12.3, 1.5/18.7, 2/25.1, 2.5/41.8, 3/68.4} {
                    \fill[white, draw=chartline, line width=5pt] (\x,\y) circle (12pt);
                }
                
                \draw[accentmain, dashed, line width=4pt] (2.4, 0) -- (2.4, 80);
                \node[above, font=\Large\bfseries, color=accentmain, fill=white, inner sep=5pt] at (2.4, 80) {2019 Crisis};
            \end{tikzpicture}
            
            \vspace{2cm}
            \begin{flushleft}
            \textit{Figure 1: All 5 CEs show accelerating convergence post-2019. The trajectory shifts from linear adoption to exponential integration.}
            \end{flushleft}
        \end{tcolorbox}

        \vspace{2cm}

        % --- 5. DISCUSSION ---
        \begin{tcolorbox}[posterbox={5}{Discussion}]
            \textbf{Why Loyalty Signaling Increases:}
            
            \vspace{1cm}
            \begin{itemize}
                \item \textbf{Crisis Response:} Political shocks (e.g., 2019 protests, NSL implementation) $\rightarrow$ increased rhetorical alignment as a reactive institutional mechanism.
                \item \textbf{Credible Signal:} Using Beijing's precise vocabulary demonstrates institutional integration and compliance to the center.
                \item \textbf{Bureaucratic Logic:} State-speak signals competence and alignment within the broader hierarchical Chinese governance system.
            \end{itemize}
        \end{tcolorbox}

    \end{column}

    % ================= RIGHT COLUMN =================
    \begin{column}{0.45\textwidth}

        % --- 2. METHODS ---
        \begin{tcolorbox}[posterbox={2}{Methods}]
            \textbf{Data \& Approach:}
            \vspace{0.5cm}
            \begin{itemize}
                \item \textbf{Corpus:} Natural language processing (NLP) of $N=892$ Policy Addresses and major speeches spanning from 1997 to 2024.
                \item \textbf{Languages:} Analyzed Traditional Chinese documents with disambiguation via \textit{Jieba} segmentation.
            \end{itemize}
            
            \vspace{1.5cm}
            \textbf{Analysis Pipeline:}
            \vspace{0.8cm}
            
            \begin{center}
            \begin{tikzpicture}[x=6cm, y=0cm]
                \tikzstyle{block} = [rectangle, fill=accentmain!20, text=textdark, align=center, rounded corners=15pt, minimum height=3cm, text width=5cm, font=\Large]
                \tikzstyle{arrow} = [thick,->,>=stealth, line width=4pt, accentmain]

                \node[block] (A) at (0,0) {\textbf{Corpus}\\[0.2em] \normalsize 892 Speeches};
                \node[block] (B) at (1.2,0) {\textbf{Dictionary}\\[0.2em] \normalsize 400+ Keywords};
                \node[block] (C) at (2.4,0) {\textbf{Density Cal.}\\[0.2em] \normalsize Score/10k chr};
                
                \draw[arrow] (A) -- (B);
                \draw[arrow] (B) -- (C);
            \end{tikzpicture}
            \end{center}
        \end{tcolorbox}

        \vspace{2cm}

        % --- 4. RESULTS 2: COMPARATIVE ---
        \begin{tcolorbox}[posterbox={4}{Comparative Analysis}]
            Cumulative \textit{guanqiang} adoption by administration. The rhetoric alignment acts as a leading indicator of actual policy convergence.\\[1.5cm]
            
            \renewcommand{\arraystretch}{2} 
            \begingroup 
            \Large
            \begin{tabular}{@{} l c r @{}}
                \toprule[4pt]
                \textbf{\color{accentmain} CE Administration} & \textbf{\color{accentmain} Density} & \textbf{\color{accentmain} Growth Rate} \\
                \midrule[2pt]
                Tung Chee-hwa (1997--2005) & 12.3 & — \\
                Donald Tsang (2005--2012) & 18.7 & \textbf{\textcolor{green!60!black}{+52\%}} \\
                C.Y. Leung (2012--2017) & 25.1 & \textbf{\textcolor{green!60!black}{+34\%}} \\
                Carrie Lam (2017--2022) & 41.8 & \textbf{\textcolor{green!60!black}{+67\%}} \\
                \rowcolor{accentmain!10}
                \textbf{John Lee (2022--2024)} & \textbf{68.4} & \textbf{\textcolor{green!60!black}{+64\%}} \\
                \bottomrule[4pt]
            \end{tabular}
            \endgroup
            
            \vspace{2cm}
            \begin{tcolorbox}[colback=white, colframe=accentline, boxrule=2pt, arc=10pt, left=1cm, right=1cm, top=1cm, bottom=1cm]
                \centering \textbf{Key Finding:} We observe a staggering \textbf{$5.6\times$ aggregate increase} in mainland political terminology across three decades.
            \end{tcolorbox}
        \end{tcolorbox}

        \vspace{2cm}

        % --- 6. CONCLUSION & REFERENCES ---
        \begin{tcolorbox}[posterbox={6}{Conclusion \& References}]
            \textbf{Conclusions:}
            \vspace{0.8cm}
            \begin{itemize}
                \item Mainlandization in Hong Kong is \textbf{systematic, not random}.
                \item Rhetorical convergence serves as political insurance and accelerates sharply during moments of central-local friction.
                \item This linguistic shift reflects a broader, substantive transition from the legacy "Two Systems" model toward a securitized "One Country" framework.
            \end{itemize}
            
            \vspace{2cm}
            \textbf{Selected References:}
            \vspace{0.5cm}
            \\ \normalsize 
            1. Smith, J. (2023). \textit{Lexical Choices as Politics}. Journal of Political Discourse.\\
            2. Wang, H. (2024). \textit{Center-Local Relations in Hong Kong}. Asian Survey.\\
            3. Chan, A. (2025). \textit{Loyalty Signaling in Bureaucracies}. Governance.
        \end{tcolorbox}

    \end{column}

\end{columns}

\end{frame}
\end{document}